\documentclass{beamer}

% Packages
\usepackage[utf8]{inputenc}
\usepackage{graphicx}
\usepackage{booktabs}
\usepackage{tikz}
\usepackage{hyperref}

%----------------------------------------------------------------------------------------
%	UTSA THEME CONFIGURATION
%----------------------------------------------------------------------------------------

% Define official UTSA Colors
% Source: UTSA Brand Identity Guide
\definecolor{utsaOrange}{HTML}{F15A22} % PMS 1665
\definecolor{utsaNavy}{HTML}{0C2340}   % PMS 289
\definecolor{utsaBlue}{HTML}{002A5C}   % Alternate/Dark Blue

% Use a clean base theme
\usetheme{Madrid}

% Custom Color Palette
\setbeamercolor{structure}{fg=utsaOrange}
\setbeamercolor{title}{bg=utsaNavy, fg=white}
\setbeamercolor{frametitle}{bg=utsaNavy, fg=white}
\setbeamercolor{section in head/foot}{bg=utsaOrange, fg=white}
\setbeamercolor{author in head/foot}{bg=utsaNavy, fg=white}
\setbeamercolor{date in head/foot}{bg=utsaNavy, fg=white}
\setbeamercolor{block title}{bg=utsaOrange, fg=white}
\setbeamercolor{block body}{bg=utsaOrange!10, fg=black}

% Add the UTSA logo to the header or footer (optional)
% \logo{\includegraphics[height=1cm]{utsa_logo.png}} % Uncomment if you have a logo file

%----------------------------------------------------------------------------------------
%	TITLE PAGE INFO
%----------------------------------------------------------------------------------------

\title[Short Title for Footer]{Complexity of Deep Computations via Topology}
\subtitle{A Model-Theoretic Approach to Neural Limits}
\author[Author Name]{Your Name Here}
\institute[UTSA]{
    Department of Mathematics\\
    The University of Texas at San Antonio\\
    \medskip
    \textit{email@utsa.edu}
}
\date{\today}

%----------------------------------------------------------------------------------------
%	PRESENTATION SLIDES
%----------------------------------------------------------------------------------------

\begin{document}

% Title Slide
\begin{frame}
    \titlepage
\end{frame}

% Outline Slide
\begin{frame}{Outline}
    \tableofcontents
\end{frame}

%------------------------------------------------
\section{Introduction}
%------------------------------------------------

\begin{frame}{Introduction}
    \begin{itemize}
        \item Motivation: Understanding the limits of deep learning models.
        \item Neural Networks as iterated functions.
        \item What happens when depth $\to \infty$?
        \item Key Question: \textbf{How complex are the limit functions?}
    \end{itemize}
\end{frame}

\begin{frame}{Key Concepts}
    \begin{block}{Deep Computations}
        Limits of iterative processes, modeled through the lens of function space topology.
    \end{block}
    
    \begin{exampleblock}{Example: Newton's Method}
        Iterating a rational map $z \mapsto z - p(z)/p'(z)$ can lead to complex fractal behavior (Julia sets) or simple convergence basins.
    \end{exampleblock}
\end{frame}

%------------------------------------------------
\section{Topological Framework}
%------------------------------------------------

\begin{frame}{Rosenthal Compacta}
    \begin{itemize}
        \item We use \textbf{Rosenthal Compacta} to classify these limit spaces.
        \item Definition: Compact subsets of the first Baire class $B_1(X)$.
        \item Connection to Model Theory:
        \begin{itemize}
            \item \textbf{NIP} (No Independence Property) $\leftrightarrow$ Tame topology.
            \item \textbf{IP} (Independence Property) $\leftrightarrow$ Wild/Chaotic topology.
        \end{itemize}
    \end{itemize}
\end{frame}

\begin{frame}{Todorčević’s Trichotomy}
    A classification of dynamical complexity:
    \begin{enumerate}
        \item \textbf{Metrizable}: The simplest tame case.
        \item \textbf{Hereditarily Separable}: Non-metrizable but manageable.
        \item \textbf{First Countable}: More complex.
        \item \textbf{General Rosenthal Compacta}: The "wilder" tame class.
    \end{enumerate}
\end{frame}

%------------------------------------------------
\section{Results}
%------------------------------------------------

\begin{frame}{Main Theorem}
    \begin{theorem}
        Let $\Delta$ be a set of computations. The space of deep computations $\overline{\Delta}$ is a Rosenthal compactum if and only if $\Delta$ satisfies the NIP.
    \end{theorem}
    
    \vspace{0.5cm}
    
    \textbf{Implications:}
    \begin{itemize}
        \item We can predict the learnability of a model based on its topological class.
        \item Provides a "Rosetta Stone" connecting Topology, Logic, and Machine Learning.
    \end{itemize}
\end{frame}

%------------------------------------------------
\section{Conclusion}
%------------------------------------------------

\begin{frame}{Summary \& Future Work}
    \textbf{Summary:}
    \begin{itemize}
        \item Established a correspondence between NIP and Rosenthal Compacta in Deep Learning.
        \item Classified limits of neural networks into 4 complexity classes.
    \end{itemize}
    
    \vspace{0.5cm}
    
    \textbf{Future Directions:}
    \begin{itemize}
        \item Exploring randomized versions of NIP.
        \item Application to specific architectures (Transformers, RNNs).
    \end{itemize}
\end{frame}

\begin{frame}{Q\&A}
    \centering
    \Large \textbf{Thank You!}
    
    \vspace{1cm}
    \footnotesize
    \textit{The University of Texas at San Antonio}
\end{frame}

\end{document}
